% ===========================================================================
% SOURCES: docs/claude_project/02_SOFTWARE_AND_QUANTIZATION.md
%          docs/phase2_export_explained.md
%          docs/references/ReferenceModelExplained.md
%          software/export/manifest.json, scripts/analyze_workload.py
%          docs/design_decisions.md: DD-009, DD-010, DD-011, DD-019
%
% START HERE. This is the first chapter to write: the work is finished, it has
% not changed since 2026-07-24, and it is documented better than any other
% part of the project.
% ===========================================================================
\chapter{Software Reference Model and Workload Analysis}
\label{chap:software}

\section{Dataset}
\label{sec:dataset}
% The four classes: \thClassList. Sources, counts, the split.

\subsection{Data Leakage and Leaf-Level Resplit}
\label{sec:leakage}
% Tell this straight: the first split reported 100.00% float and 100.00% int8,
% and those numbers were wrong because augmented duplicates of the same leaf
% appeared in both train and validation. scripts/resplit_dataset.py resplits by
% leaf. The honest accuracy on grape_clean is lower and it is the one reported
% in Section~\ref{sec:accuracy}.
%
% This belongs in the thesis rather than being quietly dropped: it is a
% failure mode endemic to the plant-disease literature (Section~\ref{sec:related-plant})
% and finding it in your own work is the evidence that the measurement
% discipline is real.

\section{Model Training}
\label{sec:training}
% Transfer learning from ImageNet, the schedule, the augmentation, the
% checkpoint. Record the checkpoint hash: \thCkptSha{} (truncated; the manifest
% carries all 64 hex digits). Every result in this thesis is against that one
% checkpoint, and saying so is what makes the rest reproducible.

\section{Post-Training Quantization to int8}
\label{sec:ptq}

\subsection{Quantization Scheme}
\label{sec:quant-scheme}
% Per-channel weights, per-tensor activations; the scales in the manifest; the
% m0/shift fixed-point requantiser. Use the notation from Section~\ref{sec:uniform-quant}.

\subsection{Residual Addition}
\label{sec:residual-add}
% \DD{009}. Two inputs at different scales into one output scale. Include the
% property gen_program.py asserts on every run (residual M <= 1) and say why a
% property checked by a script beats a property checked by a reviewer.

\subsection{Global Average Pooling}
\label{sec:gap}
% \DD{010}. The integer mean and its rounding - exact, because
% Chapter~\ref{chap:verification} compares it element by element.

\subsection{Integer Logits and Argmax}
\label{sec:logits}
% \DD{011}. The int16 classifier tail; why argmax on integer logits is
% identical to argmax on the dequantised ones.

\section{Parameter Export and Golden Vectors}
\label{sec:export}

\subsection{Export Format and Manifest}
\label{sec:manifest}
% What export.py emits: int8 weights, int32 bias, m0/shift per channel, and a
% manifest of \thManifestOps{} ops carrying ckpt_sha256. Note the guard:
% gen_program.py refuses a hash mismatch AND a missing hash. The second case is
% the one that actually bit (Section~\ref{sec:acc-bound}).

\subsection{Golden Vectors}
\label{sec:golden}
% The central artefact of the thesis. The exporter emits the exact integer
% output of every layer; the hardware is required to reproduce them bit for
% bit. State the contract here in one sentence and let Chapter~\ref{chap:verification}
% build the method on it.

\subsection{Provable Accumulator Bound}
\label{sec:acc-bound}
% THE \DD{019} STORY. Structure it as: what was done, why it was reasonable,
% why it was wrong, what replaced it.
%
%   ACC_W = 21 was justified by a MEASUREMENT: max|acc+bias| = 20,888 against
%   a limit of +-1,048,576, a 50x margin. Five weeks later the model was
%   retrained and one channel of features.3.conv.0.0 acquired a bias of
%   1,184,089 - which does not fit a 21-bit signed accumulator on its own, for
%   ANY input. The hardware wrapped negative, ReLU6 clamped to zero, and
%   2,215,904 of \thElements elements were wrong downstream, silently.
%
%   The fix was not ACC_W = 26. It was replacing the measurement with a bound
%   the exporter recomputes on every run and exits non-zero on:
%
\begin{equation}
  \lvert \mathrm{acc} + b \rvert \;\le\;
  \Bigl( \sum_{\mathrm{taps}} \lvert w_{oc,\mathrm{tap}} \rvert \Bigr)
  \cdot \max \lvert a \rvert \;+\; \lvert b_{oc} \rvert
  \label{eq:acc-bound}
\end{equation}
%
%   max|a| is a property of the number format, not of the data, so no image can
%   exceed it. Land the general lesson, because it recurs in Chapter~\ref{chap:timing}
%   and Section~\ref{sec:risks}: a retrain invalidates measurements, not bounds.

\section{Bit-Exact C Reference Model}
\label{sec:c-model}
% Op by op, the same integer arithmetic as the exporter, in C. It is the
% executable specification the RTL is written against, and it is also what
% runs on the PS in Chapter~\ref{chap:integration}.

\section{Input Preprocessing}
\label{sec:preproc}
% \DD{033}. NEW SECTION - not in the original outline, and it closes a real
% hole: the \thElements-element bit-exactness started at 000_input.hex, so
% resize and quantisation were outside every check in the project until this
% was written. Cover the resize, the normalisation, the rounding mode
% (and why fenv.h rather than a C99 version test), and the bit-exactness
% against the golden input tensor.
% Verified on the host without a board - see firmware/test/test_preproc.c.

\section{Workload Analysis}
\label{sec:workload}

\subsection{MAC Distribution per Layer Type}
\label{sec:mac-dist}
% \thMacsConv{} MACs in the convolution layers. The distribution is the whole
% argument for the architecture in Chapter~\ref{chap:architecture}:
%   pointwise \thPwMacPct\% | depthwise \thDwMacPct\% | stem \thStemMacPct\%
% Say which total you are quoting and why - \thMacs{} includes the GAP
% accumulations, which are additions without a multiply.
% Generate the table and the figure from scripts/analyze_workload.py.

\subsection{Roofline Analysis}
\label{sec:roofline-analysis}
% The model from Section~\ref{sec:roofline}, applied. Weights are \thWeightMB{} MB and
% activations stay on chip, so the design is compute-bound - which is what
% licenses spending \thDspPctAccelAxi\% of the DSPs on one array.

\section{Identifying the Operations to Accelerate}
\label{sec:what-to-accelerate}
% The conclusion the chapter has been building to, and the bridge into
% Chapter~\ref{chap:architecture}: pointwise first and biggest, depthwise separately
% because its dataflow differs, the stem once, everything else shared.
